% --- Archon physics formalization source begin ---
% archon:physics
% archon:chemistry
% archon:covers IChO2026Problems/problem_icho_2026_t4_a7.lean
% archon:source-report reports/icho_2026/problem_icho_2026_t4_a7.source.json
% archon:problem-id icho_2026_t4
% archon:part-id T4-A7
% archon:previous-part icho_2026_t4_a6
% archon:previous-part-policy natural_language_prerequisite_only

\chapter{Icho Chemistry problem icho\_2026\_t4\_a7}
\label{ch:IChO2026Problems_problem_icho_2026_t4_a7}

\paragraph{Problem source.}
\#\# Source
58th International Chemistry Olympiad, Tashkent, Uzbekistan, 2026.
Official-materials index: https://www.icho2026.uz/problems
English problem PDF: ../icho\_2026\_source/raw/theory\_problem.pdf
English solution/rubric PDF: ../icho\_2026\_source/raw/theory\_solution.pdf
Problem PDF page 39; printed local question page 3; solution starts on PDF page 36.
The page PNG is primary visual evidence for formulas, structures, tables, and figures.

\#\# T4. The Nuclear Past of Uzbekistan
T4. The nuclear past of Uzbekistan 6\% Natural uranium consists primarily of two isotopes, 235U and 238U . 235U is responsible for sustaining nuclear fission reactions. 4.1 Calculate the atomic abundance of 235U in natural uranium, assuming it consists only of isotopes 235U (235.04 a.u.) and 238U (238.05 a.u.). The main reaction occurring in nuclear power plants is the reaction where 235U absorbs a neutron and undergoes fission, releasing energy and producing three additional neutrons, thus sustaining a chain reaction. 235 92 U +1 0 n →… + 31 0n Neutrons that induce the chain reaction are produced by several ways: a) Reaction of 11B with 𝛼-particles inside the reactor core b) Interaction of 𝛾rays with 2D nuclei present in the reactor coolant water c) Reaction of 9Be with 𝛼-particles produced in an externally installed neutron source. 4.2 Write the nuclear reaction equations for (a), (b), and (c). The graph above displays the percentage yield, P(A), of fission products by mass number, A, resulting from 235U fission. 4.3 Write the most common nuclear fission equation for splitting of 235U upon neutron absorption. The pair of elements formed in this reaction are in the same group of Periodic Table. 4.4 Calculate the energy released in this reaction (ΔE, MeV), if the binding energy 𝐵𝐸(235U) = 7.59 MeV/nucleon and average binding energy for fission products 𝐵𝐸(fis.) = 8.45 MeV/nucleon. Neglect the binding energy of free neutrons. Fission neutrons are produced with an average energy of 2 MeV and slow down through many scatterings with nuclei. After several collisions, their speed decreases until their average kinetic energy matches that of the surrounding atoms or molecules. The logarithmic energy decrement, 𝜉, shows the average decrease per collision in the logarithm of the neutron energy: 𝜉= ln ( 𝐸𝑖𝑛𝑖𝑡𝑖𝑎𝑙 𝐸𝑓𝑖𝑛𝑎𝑙) 𝐸𝑖𝑛𝑖𝑡𝑖𝑎𝑙and 𝐸𝑓𝑖𝑛𝑎𝑙are the initial and final neutron energies per collision. 𝜉is a constant for each type of material and does not depend on the initial neutron energy. 4.5 Determine the average number of collisions required, n𝑐, to slow a neutron from 2 MeV to 0.012 eV, using water as the moderator. 𝜉(water)=0.948 In 1966, the Urtabulak gas field suffered an uncontrolled natural gas well blowout, releasing enormous volumes of burning methane. 4.6 Calculate Δ𝑟𝐻∘ 298 (kJ mol−1) for one mole of methane combustion reaction at 298 K. Assume all species are gaseous. Use the thermodynamic data below Δ𝑓𝐻∘ 298(CH4) -74.8 kJ mol−1 Δ𝑓𝐻∘ 298(H2O, gas) -241.8 kJ mol−1 Δ𝑓𝐻∘ 298(CO2) -393.5 kJ mol−1 𝐶𝑃(CH4) 35 J mol−1 K−1 𝐶𝑃(H2O, gas) 34 J mol−1 K−1 𝐶𝑃(O2) 29 J mol−1 K−1 𝐶𝑃(CO2) 37 J mol−1 K−1 If you did not get an answer for 4.6, use Δ𝐻298 = −750 𝑘𝐽𝑚𝑜𝑙−1 for further calculations.

\#\# Current subquestion 4.7
Calculate Δ𝑟𝐻2000 (kJ mol−1) per mole of methane combustion reaction at 2000 K. Assume all species are gaseous.

\#\# Dataset metadata
IChO 2026; T4; raw rubric points=2; kind=theory; formalization\_ready=true.

\paragraph{Shared problem context.}
T4. The nuclear past of Uzbekistan 6\% Natural uranium consists primarily of two isotopes, 235U and 238U . 235U is responsible for sustaining nuclear fission reactions. 4.1 Calculate the atomic abundance of 235U in natural uranium, assuming it consists only of isotopes 235U (235.04 a.u.) and 238U (238.05 a.u.). The main reaction occurring in nuclear power plants is the reaction where 235U absorbs a neutron and undergoes fission, releasing energy and producing three additional neutrons, thus sustaining a chain reaction. 235 92 U +1 0 n →… + 31 0n Neutrons that induce the chain reaction are produced by several ways: a) Reaction of 11B with 𝛼-particles inside the reactor core b) Interaction of 𝛾rays with 2D nuclei present in the reactor coolant water c) Reaction of 9Be with 𝛼-particles produced in an externally installed neutron source. 4.2 Write the nuclear reaction equations for (a), (b), and (c). The graph above displays the percentage yield, P(A), of fission products by mass number, A, resulting from 235U fission. 4.3 Write the most common nuclear fission equation for splitting of 235U upon neutron absorption. The pair of elements formed in this reaction are in the same group of Periodic Table. 4.4 Calculate the energy released in this reaction (ΔE, MeV), if the binding energy 𝐵𝐸(235U) = 7.59 MeV/nucleon and average binding energy for fission products 𝐵𝐸(fis.) = 8.45 MeV/nucleon. Neglect the binding energy of free neutrons. Fission neutrons are produced with an average energy of 2 MeV and slow down through many scatterings with nuclei. After several collisions, their speed decreases until their average kinetic energy matches that of the surrounding atoms or molecules. The logarithmic energy decrement, 𝜉, shows the average decrease per collision in the logarithm of the neutron energy: 𝜉= ln ( 𝐸𝑖𝑛𝑖𝑡𝑖𝑎𝑙 𝐸𝑓𝑖𝑛𝑎𝑙) 𝐸𝑖𝑛𝑖𝑡𝑖𝑎𝑙and 𝐸𝑓𝑖𝑛𝑎𝑙are the initial and final neutron energies per collision. 𝜉is a constant for each type of material and does not depend on the initial neutron energy. 4.5 Determine the average number of collisions required, n𝑐, to slow a neutron from 2 MeV to 0.012 eV, using water as the moderator. 𝜉(water)=0.948 In 1966, the Urtabulak gas field suffered an uncontrolled natural gas well blowout, releasing enormous volumes of burning methane. 4.6 Calculate Δ𝑟𝐻∘ 298 (kJ mol−1) for one mole of methane combustion reaction at 298 K. Assume all species are gaseous. Use the thermodynamic data below Δ𝑓𝐻∘ 298(CH4) -74.8 kJ mol−1 Δ𝑓𝐻∘ 298(H2O, gas) -241.8 kJ mol−1 Δ𝑓𝐻∘ 298(CO2) -393.5 kJ mol−1 𝐶𝑃(CH4) 35 J mol−1 K−1 𝐶𝑃(H2O, gas) 34 J mol−1 K−1 𝐶𝑃(O2) 29 J mol−1 K−1 𝐶𝑃(CO2) 37 J mol−1 K−1 If you did not get an answer for 4.6, use Δ𝐻298 = −750 𝑘𝐽𝑚𝑜𝑙−1 for further calculations.

\paragraph{Current subquestion.}
Calculate Δ𝑟𝐻2000 (kJ mol−1) per mole of methane combustion reaction at 2000 K. Assume all species are gaseous.

\paragraph{Recorded answer/context.}
Assume constant heat capacities. Δ𝐻(𝑇) = Δ𝐻(298 𝐾) + Δ𝐶𝑝(𝑇−298 𝐾) Δ𝐶𝑝= [𝐶𝑝(CO2) + 2𝐶𝑝(H2O)] −[𝐶𝑝(CH4) + 2𝐶𝑝(O2)] Δ𝐶𝑝= [37 + 2(34)] −[35 + 2(29)] = 12 J mol−1 K−1 Δ𝑟𝐻∘ 2000 = −802.3 kJ mol−1 + (12Jmol−1K−1)(1702 𝐾) 1000 = −781.876 kJ mol−1 Δ𝑟𝐻∘ 2000 = −781.9 kJ mol−1 For a given value of Δ𝑟𝐻∘ 298 = −750 kJ/mol alternative answer will be: Δ𝑟𝐻2000 = −781.9 𝑘𝐽𝑚𝑜𝑙−1 Δ𝑟𝐻∘ 2000 = −750 kJ mol−1 + (12 J mol−1K−1)(1702 K) 1000 = −729.576 kJ mol−1 Δ𝑟𝐻2000 = −729.6 kJ mol−1 2 points total: 1 point for the correct heat capacity change calculation. 1 point for the correct enthalpy calculation at 2000 K. 2.0 pt

\paragraph{Figure/image paths.}
\begin{itemize}
\item ../icho\_2026\_source/image/T4\_page-3.png
\item ../icho\_2026\_source/image/T4\_page-2.png
\end{itemize}

\paragraph{Reusable previous-part conclusions.}
\begin{itemize}
\item Source T4-A6. Question: Calculate Δ𝑟𝐻∘ 298 (kJ mol−1) for one mole of methane combustion reaction at 298 K. Assume all species are gaseous. Use the thermodynamic data below Reusable conclusions: Balanced reaction: CH4(𝑔) + 2O2(𝑔) → CO2(𝑔) + 2H2O(𝑔) Δ𝑟𝐻∘ 298 = [Δ𝑓𝐻∘(CO2) + 2Δ𝑓𝐻∘(H2O)] −[Δ𝑓𝐻∘(CH4) + 2Δ𝑓𝐻∘(O2)] Δ𝑟𝐻∘ 298 = ([−393.5 + 2(−241.8)] −[−74.8 + 0])kJ mol−1 = −802.3 kJ mol−1 Molar enthalpy of methane combustion at 298 K, Δ𝑟𝐻298 = −802.3 kJ mol−1 2 points for the correct answer; 1 point if the answer is incorrect but the calculation of difference in enthalpies of products and reagents exists. 2.0 pt Policy: natural\_language\_prerequisite\_only; the current target and later answers must not be assumed
\end{itemize}

\paragraph{Formalization target.}
create a compiling Lean file with sorry bodies at `IChO2026Problems/problem\_icho\_2026\_t4\_a7.lean`.
The Lean declarations must preserve every mathematical object, quantifier, hypothesis, side condition, approximation/error guarantee, and requested conclusion from the source statement.
Use Mathlib/Physlib/Chemistry names found through LeanExplore where available. Prefer the configured benchmark Base library; introduce a local definition only when the source genuinely requires an object that the environment does not expose.

\begin{definition}\leanok
[Combustion-species carrier]
\label{def:physics:icho_2026_t4_a7:combustion_species}
\lean{IChO2026Problems.CombustionSpecies}
The species carrier has exactly methane, dioxygen, carbon dioxide, and water.
\end{definition}
\begin{proof}\leanok
These are the four species in the displayed combustion equation.\end{proof}

\begin{definition}\leanok
[Element carrier]
\label{def:physics:icho_2026_t4_a7:chemical_element}
\lean{IChO2026Problems.ChemicalElement}
The atom-balance carrier consists of carbon, hydrogen, and oxygen.
\end{definition}
\begin{proof}\leanok
They are exactly the elements occurring in the four formulas.\end{proof}

\begin{definition}\leanok
[Physical phase]
\label{def:physics:icho_2026_t4_a7:physical_phase}
\lean{IChO2026Problems.PhysicalPhase}
The phase carrier distinguishes gas, liquid, solid, and aqueous forms.
\end{definition}
\begin{proof}\leanok
The gas case is required by both T4-A6 and T4-A7.\end{proof}

\begin{definition}\leanok
[Chemical entity]
\label{def:physics:icho_2026_t4_a7:chemical_entity}
\lean{IChO2026Problems.ChemicalEntity}
\uses{def:physics:icho_2026_t4_a7:combustion_species, def:physics:icho_2026_t4_a7:physical_phase}
A chemical entity pairs one combustion species with one physical phase.
\end{definition}
\begin{proof}\leanok
This preserves phase identity rather than treating it as prose.\end{proof}

\begin{definition}\leanok
[Gaseous methane]
\label{def:physics:icho_2026_t4_a7:methane_gas}
\lean{IChO2026Problems.methaneGas}
\uses{def:physics:icho_2026_t4_a7:chemical_entity}
This is methane paired with the gas phase.
\end{definition}
\begin{proof}\leanok
It names the gaseous reactant in the source equation.\end{proof}

\begin{definition}\leanok
[Gaseous dioxygen]
\label{def:physics:icho_2026_t4_a7:dioxygen_gas}
\lean{IChO2026Problems.dioxygenGas}
\uses{def:physics:icho_2026_t4_a7:chemical_entity}
This is dioxygen paired with the gas phase.
\end{definition}
\begin{proof}\leanok
It names the gaseous oxidant in the source equation.\end{proof}

\begin{definition}\leanok
[Gaseous carbon dioxide]
\label{def:physics:icho_2026_t4_a7:carbon_dioxide_gas}
\lean{IChO2026Problems.carbonDioxideGas}
\uses{def:physics:icho_2026_t4_a7:chemical_entity}
This is carbon dioxide paired with the gas phase.
\end{definition}
\begin{proof}\leanok
It names the gaseous carbon-containing product.\end{proof}

\begin{definition}\leanok
[Gaseous water]
\label{def:physics:icho_2026_t4_a7:water_gas}
\lean{IChO2026Problems.waterGas}
\uses{def:physics:icho_2026_t4_a7:chemical_entity}
This is water paired with the gas phase.
\end{definition}
\begin{proof}\leanok
It names the gaseous hydrogen-containing product.\end{proof}

\begin{definition}\leanok
[Stoichiometric coefficient]
\label{def:physics:icho_2026_t4_a7:stoichiometric_coefficient}
\lean{IChO2026Problems.StoichiometricCoefficient}
A stoichiometric coefficient is a real molar amount per mole of reaction.
\end{definition}
\begin{proof}\leanok
This permits the coefficient arithmetic used in the solution.\end{proof}

\begin{definition}\leanok
[Chemical reaction]
\label{def:physics:icho_2026_t4_a7:chemical_reaction}
\lean{IChO2026Problems.ChemicalReaction}
\uses{def:physics:icho_2026_t4_a7:chemical_entity, def:physics:icho_2026_t4_a7:stoichiometric_coefficient}
A reaction provides reactant and product coefficients for each chemical entity.
\end{definition}
\begin{proof}\leanok
It is the carrier through which the stored reaction determines
all thermodynamic stoichiometric sums.\end{proof}

\begin{definition}\leanok
[Atomic composition]
\label{def:physics:icho_2026_t4_a7:atom_count}
\lean{IChO2026Problems.atomCount}
\uses{def:physics:icho_2026_t4_a7:combustion_species, def:physics:icho_2026_t4_a7:chemical_element}
The atomic compositions are CH$_4$, O$_2$, CO$_2$, and H$_2$O, with zero for
every absent element.
\end{definition}
\begin{proof}\leanok
Reading each molecular formula gives the displayed natural-number
counts.\end{proof}

\begin{definition}\leanok
[Displayed methane-combustion reaction]
\label{def:physics:icho_2026_t4_a7:methane_combustion}
\lean{IChO2026Problems.methaneCombustion}
\uses{def:physics:icho_2026_t4_a7:chemical_reaction, def:physics:icho_2026_t4_a7:methane_gas, def:physics:icho_2026_t4_a7:dioxygen_gas, def:physics:icho_2026_t4_a7:carbon_dioxide_gas, def:physics:icho_2026_t4_a7:water_gas}
The stored reaction has coefficients $1,2$ on the methane/dioxygen reactant
side and $1,2$ on the carbon-dioxide/water product side.
\end{definition}
\begin{proof}\leanok
This is the balanced all-gas reaction used in the official
calculation.\end{proof}

\begin{definition}\leanok
[Gas-only support]
\label{def:physics:icho_2026_t4_a7:uses_only_gas_phase}
\lean{IChO2026Problems.ChemicalReaction.usesOnlyGasPhase}
\uses{def:physics:icho_2026_t4_a7:chemical_reaction, def:physics:icho_2026_t4_a7:physical_phase}
A reaction uses only the gas phase when every entity with a nonzero reactant
or product coefficient has phase gas.
\end{definition}
\begin{proof}\leanok
It makes the source's phase condition a checkable predicate.\end{proof}

\begin{definition}\leanok
[Atom balance]
\label{def:physics:icho_2026_t4_a7:atom_balanced}
\lean{IChO2026Problems.ChemicalReaction.atomBalanced}
\uses{def:physics:icho_2026_t4_a7:chemical_reaction, def:physics:icho_2026_t4_a7:atom_count, def:physics:icho_2026_t4_a7:methane_gas, def:physics:icho_2026_t4_a7:dioxygen_gas, def:physics:icho_2026_t4_a7:carbon_dioxide_gas, def:physics:icho_2026_t4_a7:water_gas}
For each of carbon, hydrogen, and oxygen, the coefficient-weighted atom count
on reactants equals that on products.
\end{definition}
\begin{proof}\leanok
This is the element-by-element conservation law for the displayed
four-species reaction.\end{proof}

\begin{definition}\leanok
[Well-formed reaction]
\label{def:physics:icho_2026_t4_a7:well_formed}
\lean{IChO2026Problems.ChemicalReaction.wellFormed}
\uses{def:physics:icho_2026_t4_a7:chemical_reaction, def:physics:icho_2026_t4_a7:uses_only_gas_phase, def:physics:icho_2026_t4_a7:atom_balanced}
A reaction is well formed when all coefficients are nonnegative, its support
is gaseous, and atoms are balanced.
\end{definition}
\begin{proof}\leanok
These are the structural conditions relevant to the source.\end{proof}

\begin{lemma}\leanok
[Methane combustion is well formed]
\label{lem:physics:icho_2026_t4_a7:methane_combustion_well_formed}
\lean{IChO2026Problems.methane_combustion_wellFormed}
\uses{def:physics:icho_2026_t4_a7:methane_combustion, def:physics:icho_2026_t4_a7:well_formed}
The displayed reaction has nonnegative gaseous coefficients and conserves C,
H, and O atoms.
\end{lemma}
\begin{proof}\leanok
Direct substitution gives the balances $1=1$, $4=4$, and $4=4$;
all nonzero entities are explicitly gaseous.\end{proof}

\begin{definition}\leanok
[Temperature unit]
\label{def:physics:icho_2026_t4_a7:temperature}
\lean{IChO2026Problems.Temperature}
Temperature is represented by its numerical value in kelvin.
\end{definition}
\begin{proof}\leanok
This is the common numerical scale of the two source temperatures.\end{proof}

\begin{definition}\leanok
[Molar heat-capacity unit]
\label{def:physics:icho_2026_t4_a7:molar_heat_capacity}
\lean{IChO2026Problems.MolarHeatCapacity}
A heat capacity is represented in $\mathrm{J\,mol^{-1}\,K^{-1}}$.
\end{definition}
\begin{proof}\leanok
This is the unit printed in the supplied data table.\end{proof}

\begin{definition}\leanok
[Molar enthalpy unit]
\label{def:physics:icho_2026_t4_a7:molar_enthalpy}
\lean{IChO2026Problems.MolarEnthalpy}
A molar enthalpy is represented in $\mathrm{kJ\,mol^{-1}}$.
\end{definition}
\begin{proof}\leanok
This is the unit printed for both requested answers.\end{proof}

\begin{definition}\leanok
[Reference temperature]
\label{def:physics:icho_2026_t4_a7:reference_temperature}
\lean{IChO2026Problems.referenceTemperature}
\uses{def:physics:icho_2026_t4_a7:temperature}
The reference temperature is $298\,\mathrm K$.
\end{definition}
\begin{proof}\leanok
It is the temperature of the T4-A6 calculation.\end{proof}

\begin{definition}\leanok
[Target temperature]
\label{def:physics:icho_2026_t4_a7:target_temperature}
\lean{IChO2026Problems.targetTemperature}
\uses{def:physics:icho_2026_t4_a7:temperature}
The requested temperature is $2000\,\mathrm K$.
\end{definition}
\begin{proof}\leanok
It is the temperature named in T4-A7.\end{proof}

\begin{definition}\leanok
[Thermodynamic reaction data]
\label{def:physics:icho_2026_t4_a7:thermodynamic_data}
\lean{IChO2026Problems.MethaneCombustionThermodynamicData}
\uses{def:physics:icho_2026_t4_a7:chemical_reaction, def:physics:icho_2026_t4_a7:chemical_entity, def:physics:icho_2026_t4_a7:molar_enthalpy, def:physics:icho_2026_t4_a7:molar_heat_capacity, def:physics:icho_2026_t4_a7:temperature}
The data record contains the stored reaction, standard formation enthalpies at
298 K, temperature-dependent heat capacities, and reaction enthalpy.
\end{definition}
\begin{proof}\leanok
Each component is an explicit input to the thermodynamic
calculation, not a hidden chemical fact.\end{proof}

\begin{definition}\leanok
[Reaction formation enthalpy at 298 K]
\label{def:physics:icho_2026_t4_a7:reaction_formation_enthalpy}
\lean{IChO2026Problems.reactionFormationEnthalpyAt298}
\uses{def:physics:icho_2026_t4_a7:thermodynamic_data, def:physics:icho_2026_t4_a7:methane_gas, def:physics:icho_2026_t4_a7:dioxygen_gas, def:physics:icho_2026_t4_a7:carbon_dioxide_gas, def:physics:icho_2026_t4_a7:water_gas}
For each of the four gaseous entities, multiply its 298-K formation enthalpy
by its stored product coefficient minus its stored reactant coefficient and
sum the terms.
\end{definition}
\begin{proof}\leanok
When the stored reaction is methane combustion, this is exactly
the products-minus-reactants formula from T4-A6.\end{proof}

\begin{definition}\leanok
[Reaction heat-capacity change]
\label{def:physics:icho_2026_t4_a7:reaction_heat_capacity_change}
\lean{IChO2026Problems.reactionHeatCapacityChange}
\uses{def:physics:icho_2026_t4_a7:thermodynamic_data, def:physics:icho_2026_t4_a7:methane_gas, def:physics:icho_2026_t4_a7:dioxygen_gas, def:physics:icho_2026_t4_a7:carbon_dioxide_gas, def:physics:icho_2026_t4_a7:water_gas}
At any temperature, $\Delta C_p$ is the four-term sum of each entity's heat
capacity multiplied by its stored product coefficient minus reactant
coefficient.  Thus the value is determined by the reaction field, not by a
hard-coded methane formula.
\end{definition}
\begin{proof}\leanok
Substitution of the displayed reaction yields
$[C_p(\mathrm{CO_2})+2C_p(\mathrm{H_2O})]-[C_p(\mathrm{CH_4})+2C_p(\mathrm{O_2})]$.\end{proof}

\begin{definition}\leanok
[T4-A6 fallback enthalpy]
\label{def:physics:icho_2026_t4_a7:fallback_enthalpy_298}
\lean{IChO2026Problems.suppliedFallbackEnthalpyAt298}
\uses{def:physics:icho_2026_t4_a7:molar_enthalpy}
The conditional fallback supplied after T4-A6 is $-750\,\mathrm{kJ\,mol^{-1}}$.
\end{definition}
\begin{proof}\leanok
It is an explicit alternate input, not the asserted T4-A6 answer.\end{proof}

\begin{definition}\leanok
[Later fallback enthalpy]
\label{def:physics:icho_2026_t4_a7:fallback_enthalpy_2000}
\lean{IChO2026Problems.suppliedFallbackEnthalpyAt2000}
\uses{def:physics:icho_2026_t4_a7:molar_enthalpy}
The later problem's $-700\,\mathrm{kJ\,mol^{-1}}$ fallback is recorded only
for subsequent calculations.
\end{definition}
\begin{proof}\leanok
It is not used to establish the T4-A7 requested result.\end{proof}

\begin{lemma}\leanok
[Formation-enthalpy calculation at 298 K]
\label{lem:physics:icho_2026_t4_a7:enthalpy_at_298}
\lean{IChO2026Problems.methane_combustion_enthalpy_at_298}
\uses{def:physics:icho_2026_t4_a7:methane_combustion, def:physics:icho_2026_t4_a7:thermodynamic_data, def:physics:icho_2026_t4_a7:reference_temperature, def:physics:icho_2026_t4_a7:reaction_formation_enthalpy}
If the stored reaction is the displayed methane combustion, the four supplied
formation enthalpies and the reaction-formation-enthalpy relation imply
$\Delta_rH^\circ(298\,\mathrm K)=-802.3\,\mathrm{kJ\,mol^{-1}}$.
\end{lemma}
\begin{proof}\leanok
The stored reaction supplies coefficients $1,2,1,2$; substitute
the four formation enthalpies in the reaction-weighted sum and evaluate as in
T4-A6.\end{proof}

\begin{theorem}\leanok
[Methane-combustion enthalpy at 2000 K]
\label{thm:physics:icho_2026_t4_a7:target}
\lean{IChO2026Problems.methane_combustion_enthalpy_at_2000}
\uses{def:physics:icho_2026_t4_a7:methane_combustion, def:physics:icho_2026_t4_a7:thermodynamic_data, def:physics:icho_2026_t4_a7:reference_temperature, def:physics:icho_2026_t4_a7:target_temperature, def:physics:icho_2026_t4_a7:reaction_heat_capacity_change, lem:physics:icho_2026_t4_a7:enthalpy_at_298}
% SOURCE: IChO 2026 Theory Solutions, T4.7, Q4-5 (read from references/icho_2026_theory_ready.jsonl)
% SOURCE QUOTE: "Assume constant heat capacities. ΔH(T) = ΔH(298 K) + ΔCp(T − 298 K)" and "ΔCp = [37 + 2(34)] − [35 + 2(29)] = 12 J mol−1 K−1"
\textit{Source: IChO 2026 Theory Solutions, T4.7.}
Assuming the given heat capacities are temperature independent on the gas
phase and the stated Kirchhoff relation holds for the stored reaction, the
reaction heat-capacity change is $12\,\mathrm{J\,mol^{-1}\,K^{-1}}$ and
$\Delta_rH^\circ(2000\,\mathrm K)=-781.876\,\mathrm{kJ\,mol^{-1}}$, within
$0.05\,\mathrm{kJ\,mol^{-1}}$ of its one-decimal presentation $-781.9$.
\end{theorem}
% SOURCE QUOTE PROOF: "Δr H°2000 = −802.3 kJ mol−1 + (12Jmol−1 K−1)(1702 K) / 1000 = −781.876 kJ mol−1"
\begin{proof}

\leanok
The stored reaction fixes the signs and coefficients in $\Delta C_p$, so the
four constant capacities give $37+2(34)-35-2(29)=12$.  Kirchhoff's relation
then gives $-802.3+12(2000-298)/1000=-781.876$.  The difference from $-781.9$
is $0.024$, below the stated rounding bound.
\end{proof}

\begin{theorem}\leanok
[Conditional 2000-K fallback calculation]
\label{thm:physics:icho_2026_t4_a7:enthalpy_at_2000_fallback}
\lean{IChO2026Problems.methane_combustion_enthalpy_at_2000_from_t4_a6_fallback}
\uses{def:physics:icho_2026_t4_a7:methane_combustion, def:physics:icho_2026_t4_a7:thermodynamic_data, def:physics:icho_2026_t4_a7:reference_temperature, def:physics:icho_2026_t4_a7:target_temperature, def:physics:icho_2026_t4_a7:reaction_heat_capacity_change, def:physics:icho_2026_t4_a7:fallback_enthalpy_298}
When the explicit $-750\,\mathrm{kJ\,mol^{-1}}$ T4-A6 fallback is supplied as
the 298-K value, the same stored-reaction correction yields
$-729.576\,\mathrm{kJ\,mol^{-1}}$, within $0.05$ of $-729.6$.
\end{theorem}
\begin{proof}

\leanok
The same reaction-weighted heat-capacity calculation gives $\Delta C_p=12$.
Insert the fallback into Kirchhoff's relation:
$-750+12(2000-298)/1000=-729.576$; the rounding difference is $0.024$.
\end{proof}
% --- Archon physics formalization source end ---
